\begin{obeactivity}
\textbf{Aktivitas 3.1: Eksperimen modulo dan invers}
\begin{enumerate}
    \item Pilih dua bilangan $a$ dan $n$.
    \item Periksa apakah $\text{PBB}(a,n) = 1$.
    \item Jika ya, carilah invers modulo $a \pmod n$ dengan mencoba nilai atau
          dengan Algoritma Euclidean.
    \item Diskusikan apa yang terjadi jika $\text{PBB}(a,n) \neq 1$.
\end{enumerate}
\end{obeactivity}

\begin{obeactivity}
\textbf{Aktivitas 3.2: Algoritma Euclidean}
Tulis program yang menerima dua bilangan bulat lalu mengembalikan PBB-nya dengan
pembagian berulang.
\begin{enumerate}
    \item Uji program dengan pasangan $(48,18)$ dan $(1970,1066)$.
    \item Tampilkan setiap langkah pembagian sampai sisanya nol.
    \item Bandingkan hasilnya dengan perhitungan manual.
\end{enumerate}
\end{obeactivity}

\begin{obeactivity}
\textbf{Aktivitas 3.3: Mengukur entropi teks nyata}
Ambil satu paragraf berita berbahasa Indonesia (minimal 200 huruf), lalu kerjakan
langkah berikut dengan program sederhana.
\begin{enumerate}
    \item Buang spasi dan tanda baca, lalu hitung frekuensi setiap huruf dan
          gambarkan diagram batangnya.
    \item Hitung entropi $H(P)$ dengan Persamaan~\ref{eq:entropi}, lalu bandingkan
          hasilnya dengan entropi maksimum $\log_2 26 = 4{,}7004$ bit/karakter.
    \item Enkripsikan paragraf itu dengan Caesar cipher, lalu hitung entropi
          cipherteksnya dan gambarkan diagram batangnya.
    \item Simpulkan mengapa cipherteks yang sebarannya lebih rata lebih sulit
          dipecahkan dengan analisis frekuensi.
\end{enumerate}
\end{obeactivity}
